% ===================================================================
%  TABEL Nr. 12 — Die stasie. — Die spoorweë. — Die skepe.
%  Meme regime que les precedents. Tableau a UNE SEULE COLONNE.
%
%  DEUX NOTES, ET ELLES NE SONT PAS AU MEME ENDROIT. Celle de
%  l'apparat — les horaires different d'un pays a l'autre, on prend
%  ceux du tableau francais — ouvre le tableau ; celle en pied,
%  t12-noto-1, dit que le voyage est decrit selon la frontiere
%  d'avant-guerre, et le fac-simile la place apres le 14. Le
%  neerlandais la laisse la ; cette colonne aussi.
%
%  LE 11 ET LE 18 ONT UN RENVOI HORS DU GRAS : « pedvarmigili » (48)
%  et « sennombra » (88) sont composes en romain, renvoi compris. Le
%  releve compte les renvois, pas les gras.
%
%  ET DEUX MOTS DEMANDENT DE NE PAS SE FIER AU FRANCAIS :
%  « biblioteko » (14) n'est pas la bibliotheque mais le kiosque a
%  livres du hall, et « trotuaro » (53, 65) n'est pas le trottoir
%  mais le quai. L'afrikaans dit boekwinkel et perron.
% ===================================================================

%%K t12-apar-1 apar
\VUsaut{4.0mm}
\VUtitre{15.0pt}{60}{TABEL Nr. 12}
\VUsaut{2.4mm}
\VUpk{11.4pt}{40}{Die stasie. --- Die spoorweë. --- Die skepe.}
\VUsaut{3.4mm}
\textsc{Nota.} --- \textit{Die uurroosters wat in elke land gebruik word is verskillend,
ons gebruik dus as voorbeeld die ure van die} \VUgras{Franse tabel.}

%%K t12-01-1 p
1. --- Ek het pas deur Duitsland gereis. Voor die vertrek het ek 'n \VUgras{spoorboekie}
\textsuperscript{(15)} gekoop om die uur van vertrek van die treine te weet. Nadat ek
vasgestel het dat die gunstigste trein om nege-uur die aand uit Parys sou vertrek en om
een-uur dertien in Straatsburg sou aankom, het ek my reis voorberei, en die volgende dag
het ek my na die stasie laat bring.

%%K t12-02-1 p
2. --- Toe ek die groot vertrekhal binnegegaan het, agter 'n ou plattelandse
\VUgras{priester} \textsuperscript{(2)} wat sy \VUgras{reissak} \textsuperscript{(3)} in
die hand gedra het, het baie \VUgras{reisigers} \textsuperscript{(1)} al aangekom.

%%K t12-02-2 p
Omdat ek 'n \VUgras{telegram} \textsuperscript{(5)} wou stuur, het ek my deur 'n
\VUgras{kontroleur} \textsuperscript{(6)} \VUgras{die telegraafkantoor}
\textsuperscript{(4)} laat aanwys.

%%K t12-03-1 p
3. --- Voordat ek die \VUgras{wagkamer} \textsuperscript{(7)} binnegegaan het, het ek my
van 'n \VUgras{kaartjie} \textsuperscript{(9)} heen en terug gaan voorsien.

%%K t12-04-1 p
4. --- Ek het nie lank by die \VUgras{loket} \textsuperscript{(8)} gewag nie, want daar
was min mense. 'n \VUgras{Soldaat} \textsuperscript{(10)}, wat met verlof vertrek het,
was so hoflik om sy beurt aan my af te staan, sodat ek genoeg tyd gehad het om 'n paar
inligtings by die \VUgras{stasiemeester} \textsuperscript{(13)} te vra, wat met die
\VUgras{kommissaris} \textsuperscript{(12)} van die toesig en met die diensdoende
\VUgras{gendarme} \textsuperscript{(11)} gepraat het.

%%K t12-tit-1 sub
\VUpk{10.2pt}{40}{Die inskrywing van die bagasie.}
\VUsaut{3.0mm}

%%K t12-05-1 p
5. --- Nadat ek 'n koerant by die \VUgras{boekwinkel} \textsuperscript{(14)} gekoop het,
het ek my \VUgras{bagasie} \textsuperscript{(17)} laat weeg wat 'n \VUgras{kruier}
\textsuperscript{(16)} pas op 'n \VUgras{bagasiekar} \textsuperscript{(19)} aangevoer
het; \VUgras{die beampte} \textsuperscript{(22)} wat met die \VUgras{weegbrug}
\textsuperscript{(21)} belas is, het my laat weet dat my koffer vyf en dertig kilogram
weeg; dit het 'n \VUgras{oorgewig} van vyf kilogram gemaak, waarvoor ek 'n bybetaling
skuld, by die \VUgras{loket} van die bagasie \textsuperscript{(23)}.

%%K t12-06-1 p
6. --- Omdat ek te vroeg was, het ek vrye tyd gehad om die verkeer van die reisigers op
die stasie te bekyk. Sommige het hulle handbagasie by die \VUgras{bagasiebewaarplek}
\textsuperscript{(25)} afgegee of afgehaal; ander het die \VUgras{uurroosters}
\textsuperscript{(26)} geraadpleeg. Die aankomende reisigers het hulle na die
\VUgras{uitgang} \textsuperscript{(32)} gehaas met hulle \VUgras{koffer}
\textsuperscript{(24)} in die hand. Enkeles het by die \VUgras{bagasie-uitgifte}
\textsuperscript{(27)} gewag en hulle \VUgras{bewys} \textsuperscript{(29)} getoon om
hulle koffers, \VUgras{kiste} \textsuperscript{(28)} of \VUgras{mandjies}
\textsuperscript{(30)} terug te kry.

%%K t12-07-1 p
7. --- \VUgras{Die beamptes van die aksyns} \textsuperscript{(33)} het die pakke
sorgvuldig ondersoek om vas te stel of 'n mens iets aan te gee het.

%%K t12-08-1 p
8. --- Sommige reisigers het na die \VUgras{buffet} \textsuperscript{(34)} gegaan, waar
'n \VUgras{kelner} \textsuperscript{(35)} hom beywer het om hulle te bedien. Hulle moet
haas, want die \VUgras{trein} \textsuperscript{(36)}, wat 'n sneltrein is, sal nie lank
op die stasie bly nie.

%%K t12-09-1 p
9. --- Daarna het ek na die vertrekperron gegaan en 'n deurgaande \VUgras{wa}
\textsuperscript{(37)} gesoek om nie gedwing te wees om tydens die reis van wa te
verwissel nie. Ek het in 'n gangrytuig geklim wat 'n \VUgras{kompartement}
\textsuperscript{(38)} vir rokers bevat en 'n kompartement wat vir «dames alleen»
gereserveer is.

%%K t12-tit-2 sub
\VUpk{10.2pt}{40}{Vertrek in die nag.}
\VUsaut{3.0mm}

%%K t12-10-1 p
10. --- Nadat ek die \VUgras{trapplank} \textsuperscript{(40)} opgeklim, die
\VUgras{deurtjie} \textsuperscript{(39)} toegemaak, die \VUgras{gordyntjie}
\textsuperscript{(41)} opgerol en my handbagasie op die \VUgras{net}
\textsuperscript{(43)} gesit het, het ek op die \VUgras{bankie} \textsuperscript{(42)}
gaan sit. Toe ek die \VUgras{lampenis} \textsuperscript{(44)} die lampe sien aansteek
het, het ek gedink dat ek 'n nag in die wa sou moet deurbring en ek het die beampte
geroep wat met die verhuur van die \VUgras{hoofkussings} \textsuperscript{(45)} belas is
om een daarvan te vra.

%%K t12-11-1 p
11. --- Die kondukteur van die trein het die kaartjies kom nagaan. Ek het hom gevra
waarom ons nog nie vertrek het nie, want dit is die regte uur. Hy het my geantwoord dat
die \VUgras{treinleier} \textsuperscript{(46)} besig is om fietse in die
\VUgras{bagasiewa} \textsuperscript{(47)} te laat sit. Boonop het 'n mens nog nie al die
voetwarmers \textsuperscript{(48)} verwissel nie.

%%K t12-12-1 p
12. --- Maar ons sou gou vertrek, want die \VUgras{stoker} \textsuperscript{(50)}, op sy
\VUgras{tender} \textsuperscript{(49)}, het steenkool gevat om die vuur van die
\VUgras{lokomotief} \textsuperscript{(54)} aan te wakker, en \VUgras{die masjinis}
\textsuperscript{(51)} het net op die sein van die \VUgras{onderstasiemeester}
\textsuperscript{(52)} gewag, wat op die \VUgras{perron} \textsuperscript{(53)} was.

%%K t12-13-1 p
13. --- Die \VUgras{ketel} \textsuperscript{(56)} van die lokomotief het dof gegrom; die
\VUgras{skoorsteen} \textsuperscript{(55)} het strome rook uitgebraak vermeng met
\VUgras{stoom} \textsuperscript{(60)}. 'n Skerp \VUgras{fluit} \textsuperscript{(59)} het
weerklink en die \VUgras{suiers} \textsuperscript{(58)} het begin beweeg en die
\VUgras{wiele} \textsuperscript{(57)} van die swaar masjien aangedryf.

%%K t12-tit-3 sub
\VUsaut{3.0mm}
\VUpk{10.2pt}{40}{Vertrek!}
\VUsaut{3.0mm}

%%K t12-14-1 p
14. --- Die snelheid van die trein, wat op die \VUgras{spore} \textsuperscript{(64)}
geloop het, is versnel; die \VUgras{perrons} \textsuperscript{(65)} het verdwyn; ons het
verby die \VUgras{gemakhuisies} \textsuperscript{(66)} gegaan, en ook 'n ry waens wat op
die ander \VUgras{spoor} \textsuperscript{(63)} gerangeer staan: \VUgras{slaapwaens}
\textsuperscript{(67)}, 'n \VUgras{eetwa} \textsuperscript{(68)}, 'n \VUgras{poswa}
\textsuperscript{(69)}.

%%K t12-noto-1 noto
\VUnotes{143.2mm}{%
\textit{(*) Nota.} --- \textit{Vanselfsprekend is die beskrywing van die reis volgens die
grens van voor die oorlog.}%
}

%%K t12-15-1 p
15. --- Daarna het die \VUgras{goederestasie} \textsuperscript{(71)} voor ons oë
verbygetrek. Onder die loodse het beamptes die \VUgras{goedere} \textsuperscript{(72)}
gelaai wat met 'n stadige trein versend word. \VUgras{Rangeerders}
\textsuperscript{(76)} naby die \VUgras{waterreservoir} \textsuperscript{(73)} het
\VUgras{veewaens} \textsuperscript{(75)} en \VUgras{plat waens} \textsuperscript{(79)}
gerangeer om 'n \VUgras{goederetrein} \textsuperscript{(80)} saam te stel.

%%K t12-16-1 p
16. --- Ons het oor die laaste \VUgras{wissel} \textsuperscript{(78)} gegaan, waarby die
wisselwagter gestaan het; ons het verby die \VUgras{seinskyf} \textsuperscript{(86)}
gekom en die stasie het in die verte verdwyn.

%%K t12-17-1 p
17. --- Gou het ons oor 'n \VUgras{ysterbrug} \textsuperscript{(74)} gery wat oor die
\VUgras{rivier} \textsuperscript{(81)} loop. Nadat ons daardie brug oorgesteek het, het
ons langs die stroom gery, waarop kragtige \VUgras{sleepbote} \textsuperscript{(82)}
swaar \VUgras{vragskuite} \textsuperscript{(83)} getrek het. Ons het boonop 'n
\VUgras{passasiersskip} \textsuperscript{(84)} gesien toe dit by 'n \VUgras{ponton}
\textsuperscript{(85)} aangelê het.

%%K t12-18-1 p
18. --- Nadat ons blitsvinnig verby etlike stasies gegaan het, het ons deur 'n paar
\VUgras{tonnels} \textsuperscript{(87)} gery en agter ons ontelbare \VUgras{huisies}
\textsuperscript{(88)} van \VUgras{oorwegwagters} gelaat. Gewieg deur die skud van die
trein het ek diep aan die slaap geraak.

%%K t12-tit-4 sub
\VUsaut{3.0mm}
\VUpk{10.2pt}{40}{Stasie Deutsch-Avricourt.}
\VUsaut{3.0mm}

%%K t12-19-1 p
19. --- Ek is skielik wakker gemaak deur die stem van 'n beampte wat geroep het:
«Deutsch-Avricourt! \textsuperscript{(*)}. Almal uitklim!» Die inspeksie van die doeane
en al die vervelige formaliteite van die grens het plaasgevind. Ek moes my sakhorlosie
gelykstel met die \VUgras{groot klok} \textsuperscript{(89)} van die stasie, wat vyf en
vyftig minute voorgeloop het; skaars wakker, het ek die \VUgras{groot wyser}
\textsuperscript{(90)} moeilik van die \VUgras{klein een} \textsuperscript{(91)}
onderskei; die \VUgras{wyserplaat} \textsuperscript{(92)} self was heeltemal verward deur
die oggendmis.

%%K t12-20-1 p
20. --- Terwyl die doeanebeamptes hulle inspeksie gedoen het, het ek gapend die
\VUgras{plakkate} \textsuperscript{(93)} ontsyfer wat die stasiemure versier. Ek het
ontbyt geëet om die tyd te verdryf, die kaart van die Duitse \VUgras{net}
\textsuperscript{(94)} geraadpleeg om te sien hoeveel afstand my nog van Straatsburg
skei, en ek het my wa weer binnegegaan, versterk deur die hoop om gou die doel van my
reis te bereik.
\VUsaut{6.0mm}
\VUfilet{22mm}
